\begin{tikzpicture}[font=\scriptsize, >=latex]
\node[circle, draw, thick, fill=goodbg, minimum size=1.7cm] (run) at (0,0) {RUN};
\node[circle, draw, thick, fill=amberbg, minimum size=1.7cm] (halt) at (5.4,0) {HALTED};
\draw[->, thick] (-2.2,0) -- (run) node[midway, above]{\texttt{reset}};
\draw[->, thick, loop above] (run) edge node[align=center]{fetch--decode--execute\\one cycle per instruction} ();
\draw[->, thick] (run) -- (halt) node[midway, above]{\texttt{HALT} decoded};
\draw[->, thick] (halt) .. controls (3.4,-2.2) and (2.0,-2.2) .. (run)
  node[midway, below]{\texttt{reset}};
\node[align=center, text=slate, text width=8.6cm] at (2.7,-3.4)
{Two states are enough because the datapath is combinational: each RUN
edge completes a whole instruction. A multi-cycle design (a memory
load that takes two beats, say) would add states per instruction phase
--- same FSM pattern, more nodes.};
\end{tikzpicture}
